\documentclass[a4paper,10pt]{article}
\usepackage[utf8]{inputenc}
\usepackage[T1]{fontenc}
\usepackage[french]{babel}
\usepackage[margin=1.5cm]{geometry}
\usepackage{xcolor}
\usepackage{tcolorbox}
\usepackage{enumitem}
\usepackage{multicol}
\usepackage{lmodern}
\usepackage{tabularx}
\usepackage{booktabs}

\renewcommand{\familydefault}{\sfdefault}
\pagestyle{empty}
\setlist{nosep, leftmargin=1.5em}

\definecolor{navy}{RGB}{0,51,102}
\definecolor{teal}{RGB}{0,120,120}
\definecolor{crimson}{RGB}{180,0,0}
\definecolor{green}{RGB}{0,120,50}
\definecolor{lightbg}{RGB}{245,248,252}
\definecolor{orange}{RGB}{200,100,0}
\definecolor{brown}{RGB}{120,60,0}
\definecolor{darkred}{RGB}{140,0,0}

\tcbset{basebox/.style={
  colback=lightbg, colframe=navy, coltitle=white,
  fonttitle=\bfseries, arc=2mm, boxrule=0.5mm,
  left=3mm, right=3mm, top=2mm, bottom=2mm, after skip=6pt
}}

\newcommand{\key}[1]{\textbf{\textcolor{navy}{#1}}}
\newcommand{\warn}[1]{\textbf{\textcolor{crimson}{#1}}}
\newcommand{\en}[1]{\textit{\textcolor{teal}{#1}}}

\begin{document}

\begin{center}
  \colorbox{brown}{\parbox{\dimexpr\textwidth-2\fboxsep}{
    \centering\vspace{3pt}
    {\LARGE\bfseries\textcolor{white}{CULTURE GÉNÉRALE \& ANGLAIS — Fiche de révision}}\\
    {\normalsize\textcolor{white}{BTS SIO — Maël Mignard}}
    \vspace{3pt}
  }}
\end{center}
\vspace{6pt}

\begin{multicols}{2}

% ============================================================
\begin{tcolorbox}[basebox, colframe=brown, title=1. Culture G — Méthodologie]
\textbf{Épreuve de culture générale :}\\
Analyse d'un corpus de documents (textes, images, chiffres).\\[4pt]
\textbf{Structure de la réponse :}
\begin{enumerate}
  \item \key{Introduction} : présenter le corpus et la problématique
  \item \key{Développement} : analyse des documents (thèse, arguments, exemples)
  \item \key{Conclusion} : synthèse et ouverture
\end{enumerate}
\vspace{4pt}
\textbf{Outils d'analyse :}
\begin{itemize}
  \item Repérer la thèse défendue
  \item Identifier les arguments et exemples
  \item Décrypter le contexte (auteur, date, support)
  \item Confronter les points de vue des documents
\end{itemize}
\end{tcolorbox}

% ============================================================
\begin{tcolorbox}[basebox, colframe=brown, title=2. Culture G — Thèmes traités]
\textbf{Thèmes vus en cours :}
\begin{itemize}
  \item \key{Les inégalités} : économiques, sociales, numériques
  \item \key{La surconsommation} : société de consommation, obsolescence programmée
  \item \key{Les Temps Modernes} (Chaplin) : travail à la chaîne, déshumanisation
  \item \key{Le Truman Show} : surveillance, société du spectacle
  \item \key{Le journalisme} : liberté de la presse, fake news
  \item \key{L'histoire} : modernité, progrès, crises
\end{itemize}
\vspace{4pt}
\textbf{Notions clés :}
\begin{itemize}
  \item \key{Aliénation} : perte de liberté face au système (Marx)
  \item \key{Surveillance} : contrôle social, Big Brother (Orwell, 1984)
  \item \key{Progrès} : technique vs humain
  \item \key{Liberté d'expression} : droit fondamental (art. 11 DDHC)
\end{itemize}
\end{tcolorbox}

% ============================================================
\begin{tcolorbox}[basebox, colframe=darkred, title=3. Culture G — Auteurs et œuvres de référence]
\begin{tabular}{lll}
\toprule
\textbf{Œuvre} & \textbf{Auteur} & \textbf{Thème} \\
\midrule
1984 & Orwell & Totalitarisme, surveillance \\
Le Meilleur des Mondes & Huxley & Société contrôlée \\
Les Temps Modernes & Chaplin & Travail, déshumanisation \\
Le Truman Show & Weir & Spectacle, surveillance \\
La Société du Spectacle & Debord & Médias, consommation \\
\bottomrule
\end{tabular}
\vspace{4pt}\\
\key{Taylorisme} : organisation scientifique du travail (division des tâches) \\
\key{Fordisme} : production de masse, travail à la chaîne \\
\key{DDHC} (1789) : Déclaration des Droits de l'Homme
\end{tcolorbox}

% ============================================================
\begin{tcolorbox}[basebox, colframe=teal, title=4. ANGLAIS — Vocabulaire informatique]
\begin{tabular}{ll}
\toprule
\textbf{Anglais} & \textbf{Français} \\
\midrule
network & réseau \\
firewall & pare-feu \\
server & serveur \\
database & base de données \\
backup & sauvegarde \\
patch / update & correctif / mise à jour \\
bug & anomalie, erreur \\
software & logiciel \\
hardware & matériel \\
bandwidth & bande passante \\
latency & latence \\
uptime & temps de fonctionnement \\
downtime & temps d'arrêt \\
troubleshoot & diagnostiquer, dépanner \\
deployment & déploiement \\
workflow & flux de travail \\
stakeholder & partie prenante \\
breach & violation (de données) \\
password & mot de passe \\
encryption & chiffrement \\
\bottomrule
\end{tabular}
\end{tcolorbox}

% ============================================================
\begin{tcolorbox}[basebox, colframe=teal, title=5. ANGLAIS — Expressions professionnelles]
\textbf{Lors d'un entretien ou d'une présentation :}
\begin{itemize}
  \item \en{I'd like to introduce myself...} (se présenter)
  \item \en{My name is ... and I'm currently studying IT.}
  \item \en{I'm responsible for...} (je suis responsable de)
  \item \en{As you can see in the chart...} (comme le montre le graphique)
  \item \en{In conclusion, I would like to emphasize...}
  \item \en{Could you clarify...?} (pouvez-vous préciser)
  \item \en{To sum up / To summarize...}
\end{itemize}
\vspace{4pt}
\textbf{Pour un e-mail professionnel :}
\begin{itemize}
  \item \en{Dear Mr./Ms. [Name],}
  \item \en{I am writing to...}
  \item \en{Please find attached...} (veuillez trouver ci-joint)
  \item \en{I look forward to hearing from you.}
  \item \en{Yours sincerely / Best regards,}
\end{itemize}
\end{tcolorbox}

% ============================================================
\begin{tcolorbox}[basebox, colframe=teal, title=6. ANGLAIS — Structure d'un texte argumentatif]
\textbf{Introduction :}
\begin{itemize}
  \item \en{This document deals with...}
  \item \en{The author argues that...}
  \item \en{The main issue is...}
\end{itemize}
\vspace{4pt}
\textbf{Développement :}
\begin{itemize}
  \item \en{First of all / Firstly,...}
  \item \en{Furthermore / Moreover / In addition,...}
  \item \en{On the other hand,...} (d'un autre côté)
  \item \en{For instance / For example,...}
  \item \en{According to the author,...}
\end{itemize}
\vspace{4pt}
\textbf{Conclusion :}
\begin{itemize}
  \item \en{To conclude / In conclusion,...}
  \item \en{All things considered,...}
  \item \en{This raises the question of...}
\end{itemize}
\end{tcolorbox}

% ============================================================
\begin{tcolorbox}[basebox, colframe=navy, title=7. ANGLAIS — Temps et grammaire essentiels]
\begin{tabular}{lll}
\toprule
\textbf{Temps} & \textbf{Forme} & \textbf{Usage} \\
\midrule
Présent simple & do/does & habitudes, faits \\
Présent continu & am/is/are + -ing & action en cours \\
Prétérit & V-ed/irrég. & passé révolu \\
Present perfect & have/has + V3 & passé lié au présent \\
Futur (will) & will + V inf. & décision, prédiction \\
Conditionnel & would + V inf. & hypothèse, politesse \\
\bottomrule
\end{tabular}
\vspace{4pt}\\
\textbf{Voix passive :} \en{The server \textbf{was} configured \textbf{by} the admin.}\\
\textbf{Modaux :} can, could, must, should, may, might\\
Ex: \en{You \textbf{must} use HTTPS.} / \en{You \textbf{should} update regularly.}
\end{tcolorbox}

% ============================================================
\begin{tcolorbox}[basebox, colframe=orange, title=8. ANGLAIS — Vocabulaire CEJM en anglais]
\begin{tabular}{ll}
\toprule
\textbf{Anglais} & \textbf{Français} \\
\midrule
supply and demand & offre et demande \\
market share & part de marché \\
shareholder & actionnaire \\
CEO & PDG \\
employee / staff & salarié / personnel \\
contract & contrat \\
wage / salary & salaire \\
turnover & chiffre d'affaires \\
profit / loss & bénéfice / perte \\
supply chain & chaîne d'approvisionnement \\
merger & fusion (d'entreprises) \\
outsourcing & externalisation \\
compliance & conformité \\
regulation & réglementation \\
\bottomrule
\end{tabular}
\end{tcolorbox}

% ============================================================
\begin{tcolorbox}[basebox, colframe=green, title=9. Activités Transversales — UX/UI]
\key{UX} (User eXperience) : expérience globale de l'utilisateur \\
\key{UI} (User Interface) : interface graphique \\[4pt]
\textbf{Principes UX :}
\begin{itemize}
  \item \key{Utilisabilité} : facile à utiliser
  \item \key{Accessibilité} : utilisable par tous (handicap)
  \item \key{Cohérence} : même style partout
  \item \key{Feedback} : réponse aux actions utilisateur
  \item \key{Affordance} : l'objet suggère son usage
\end{itemize}
\vspace{4pt}
\textbf{Outils :} Figma, Adobe XD, Sketch \\
\key{Wireframe} : maquette filaire (structure) \\
\key{Prototype} : maquette interactive \\
\key{Test utilisateur} : observer de vrais utilisateurs
\end{tcolorbox}

% ============================================================
\begin{tcolorbox}[basebox, colframe=green, title=10. Activités Transversales — Veille technologique]
\key{Veille technologique} : surveiller l'évolution des technologies pour rester à jour.\\[4pt]
\textbf{Types de veille :}
\begin{itemize}
  \item \key{Veille techno} : nouvelles technologies, outils
  \item \key{Veille concurrentielle} : surveiller la concurrence
  \item \key{Veille réglementaire} : nouvelles lois et normes
  \item \key{Veille sécurité} : nouvelles failles et menaces
\end{itemize}
\vspace{4pt}
\textbf{Outils de veille :}
\begin{itemize}
  \item \key{RSS} (Really Simple Syndication) : flux d'actualités
  \item \key{Google Alerts} : alertes e-mail sur mots-clés
  \item \key{Twitter/LinkedIn} : réseaux professionnels
  \item \key{CVE} : liste des failles connues (cybersécurité)
  \item \key{ANSSI} : agence nationale de sécurité informatique
\end{itemize}
\vspace{4pt}
\key{GitLab} : plateforme DevOps (CI/CD, issues, wiki, versions)
\end{tcolorbox}

\end{multicols}

\vspace{4pt}
\begin{tcolorbox}[basebox, colframe=brown, title=Points clés à retenir pour l'examen]
\begin{multicols}{2}
\begin{itemize}
  \item Analyse de corpus : thèse + arguments + exemples dans chaque doc
  \item Introduction = présenter le corpus + poser la problématique
  \item Les Temps Modernes = Chaplin = critique du taylorisme/fordisme
  \item 1984 = Orwell = surveillance, totalitarisme, Big Brother
  \item \en{Please find attached} = veuillez trouver ci-joint
  \item Present perfect : have/has + V3 (passé lié au présent)
  \item Voix passive en anglais = be + V3
  \item UX = expérience ; UI = interface (l'UI fait partie de l'UX)
  \item Wireframe = maquette filaire ; Prototype = interactif
  \item Veille techno : RSS, Google Alerts, CVE pour les failles
  \item \warn{Freedom of speech} = liberté d'expression
  \item \en{According to} = selon, d'après (pour citer un auteur)
\end{itemize}
\end{multicols}
\end{tcolorbox}

\end{document}
